\section{Dreams as Synthesis, Not Revelation}\label{sec:deconstructing}

While compelling, these anecdotes can foster a myth of \textit{de novo} revelation. A more critical analysis suggests the dreaming brain is not a passive receiver but an active synthesizer, conducting what has been termed ``idea-incubation''.

The case of Dmitri Mendeleev's periodic table provides this essential nuance. Popular lore holds that he discovered the table in a dream. However, archival research suggests a different timeline: Mendeleev had \textit{already discovered} the periodic system \textit{before} the alleged dream took place \citep{gordin2004mendeleev}. He had been intensely preoccupied with classifying the 56 known elements, convinced a pattern existed that he ``just couldn't quite grasp''. Overcome by exhaustion, he fell asleep and, as he recounted in his diary: ``I saw in a dream a table where all the elements fell into place as required''. The dream, therefore, was not the moment of initial discovery but a process of \textit{refining the representation}—a final, non-conscious synthesis of knowledge he already possessed.

This process is not random. It is invariably preceded by intense waking obsession. Kekulé's ``extensive knowledge of organic chemistry'' and Howe's ``significant psychological and material pressure'' primed their unconscious minds \citep{findlay1965chemistry, kaempffert1924invention}. The dream state, freed from linear, conscious constraints, could then access, juxtapose, and synthesize this knowledge into a novel solution.

